\documentclass{article}
\usepackage[UTF8]{ctex}
\usepackage{amsmath,amssymb,amsfonts}
\usepackage{enumitem}
\usepackage{geometry}
\geometry{a4paper, margin=1in}
% === 错题本排版优化（不参与题目标识） ===
\setlength{\parindent}{0pt}
\setlength{\parskip}{0.35em plus 0.10em minus 0.08em}
\setlength{\emergencystretch}{3em}
\setlist{itemsep=0.25em, topsep=0.40em, parsep=0pt}
\allowdisplaybreaks[2]
% === 错题本排版优化结束 ===
\begin{document}
\section*{天津大学2024年高等代数真题}
\section*{51．天津大学2024年高等代数试卷（回忆版）}

\subsection{题目1}
1．（15 分）设 \(f(x)\) 在复数域上无重根，且

\[
f(f(x))=f^{n}(x)+a_{n-1} f^{n-1}(x)+\cdots+a_{1} f(x)+a_{0} .
\]

（1）证明：\(f(x)\) 的系数为整数．\\
（2）证明：若 \(a_{0}, a_{1}, \cdots, a_{n-1}\) 都是奇数，则 \(f(x)\) 无整数根．

\subsection{题目2}
2．（20 分）设 \(a_{1}, a_{2}, \cdots, a_{n} \neq 0\) ，记 \(m=\prod_{i=1}^{n} a_{i}\) ，且

\[
D_{n}(\lambda)=\operatorname{det}\left(\begin{array}{cccc}
1 & 1 & \cdots & 1 \\
a_{1} & a_{2} & \cdots & a_{n} \\
a_{1}^{2} & a_{2}^{2} & \cdots & a_{n}^{2} \\
\vdots & \vdots & & \vdots \\
a_{1}^{n-2} & a_{2}^{n-2} & \cdots & a_{n}^{n-2} \\
a_{1}^{n}+\frac{\lambda m}{m} & a_{2}^{n}+\frac{\lambda m}{m} & \cdots & a_{n}^{n}+\frac{\lambda m}{m}
\end{array}\right)
\]

（1）求 \(D_{n}(\lambda)\) 。\\
（2）记 \(S=\sum_{i=1}^{n} a_{i}, D_{a}(\lambda) \neq 0\) ，求

\[
D=\operatorname{det}\left(\begin{array}{cccc}
x & a_{2} & \cdots & a_{n} \\
a_{1} & x & \cdots & a_{n} \\
\vdots & \vdots & & \vdots \\
a_{1} & a_{2} & \cdots & x
\end{array}\right)+(-1)^{n}(x+s) .
\]

\subsection{题目3}
3．（15 分）设 \(\alpha>0, A=\left(\begin{array}{cc}1 & \frac{\alpha}{n} \\ \frac{\alpha}{n} & 1\end{array}\right)\) ．\\
（1）求 \(B=\lim _{n \rightarrow \infty} A^{n}\) ．\\
（2）设上一问的 \(B=\left(\begin{array}{ll}a & b \\ c & d\end{array}\right)\) ，同时

\[
D=\left(\begin{array}{lll}
a & 1 & \\
& a & \\
& & 1 \\
& & a
\end{array}\right)
\]

求 \(\operatorname{det} D\) ．

\subsection{题目4}
4．未知。

\subsection{题目5}
5．设 \(A, B\) 都是正定的实对称 \(n\) 阶方阵，证明：

\begin{itemize}
  \item[（1）] \(A+B\) 为正定的实对称方阵。
  \item[（2）] 存在 \(t>0\) ，满足 \(A-t B\) 为正定的实对称矩阵．
\end{itemize}

\subsection{题目6}
6．设 \(M_{n}(\mathbb{R})\) 为实数域上所有 \(n\) 阶矩阵构成的线性空间，记

\[
V=\left\{A \in M_{n}(\mathbb{R}) \mid \operatorname{tr}(\mathcal{A})=0\right\} .
\]

（1）证明：\(V\) 为 \(M_{\mathrm{n}}(\mathbb{R})\) 的子空间。\\
（2）求 \(V\) 的基与维数．\\
（3）定义 \(M_{\mathrm{n}}(\mathbb{R})\) 上的内积为 \((A, B)=\operatorname{tr}\left(A B^{\mathrm{T}}\right)\) ，证明：\(M_{n}(\mathbb{R})\) 按这个运算构成欧几里得空间，并求：空间 \(V\) 的正交补空间。

\subsection{题目7}
7．证明：若 \(n\) 阶矩阵 \(A\) 的迹为 0，则 \(A\) 相似于一个对角线元素都为 0 的矩阵。

\subsection{题目8}
8．存在实数 \(c_{1}, c_{2}, \cdots, c_{d}\) ，其中 \(c_{d} \neq 0\) ，

\[
a_{k+d}=c_{1} a_{k+d-1}+c_{2} a_{k+d-2}+\cdots+c_{d} a_{k} .
\]

证明：存在矩阵 \(A\) 满足 \(V_{k+1}=A V_{k}\) ，其中 \(V_{k}=\left(a_{k}, a_{k+1}, \cdots, a_{k+d-1}\right)^{\mathrm{T}}\) 的 \(d\) 维向量．\\
（1）求矩阵 \(A\) 。\\
（2）求 \(A\) 的特征多项式。\\
（3）\(A\) 的极小多项式为关于 \((x-\lambda)^{d}\) 的多项式，证明：存在次数不超过 \(d-1\) 的多项式 \(f(x)\) ，满足\(b_{k}=p(x)(x-\lambda)^{d}\) ，其中

\[
b_{k+d}=c_{1} b_{k+d-1}+c_{2} b_{k+d-2}+\cdots+c_{d} b_{k} .
\]

\subsection{题目9}
9．设 \(\alpha_{1}, \alpha_{2}, \cdots, \alpha_{n}\) 为欧氏空间 \(V\) 的一组基，\(\beta_{1}, \beta_{2}, \cdots, \beta_{n} \in V\) 满足

\[
\left(\alpha_{i}, \beta_{j}\right)= \begin{cases}1, & i=j \\ 0, & i \neq j\end{cases}
\]

\begin{itemize}
  \item[（1）] 证明：满足条件的 \(\beta_{1}, \beta_{2}, \cdots, \beta_{n}\) 是存在且唯一的．
  \item[（2）] 证明：\(\beta_{1}, \beta_{2}, \cdots, \beta_{n}\) 为 \(V\) 的一组基。
\end{itemize}
\end{document}
